\section{The \texttt{Set\_pol} McStas Component}
(Unphysical) way of setting the polarization.

\subsection*{Identification}
\begin{itemize}
  \item \textbf{Site:} 
  \item \textbf{Author:} Peter Christiansen
  \item \textbf{Origin:} Risoe
  \item \textbf{Date:} August 2006
\end{itemize}

\subsection*{Description}
\begin{lstlisting}
This component has no physical size (like Arm - also drawn that
way), but is used to set the polarisation in one of four ways:
1) (randomOn=0, normalize=0) Hardcode the polarisation to the vector (px, py, pz)
2) (randomOn!=0, normalize!=0) Set the polarisation to a random vector on the unit sphere
3) (randomOn!=0, normalize=0) Set the polarisation to a radnom vector within the unit sphere
4) (randomOn=0, normalize!=0) Hardcode the polarisation to point in the direction of (px,py,pz) but with polarization=1.
Example: Set_pol(px=0, py=-1, pz=0)
\end{lstlisting}

\subsection*{Input parameters}
Parameters in \textbf{boldface} are required; the others are optional.

\begin{longtable}{p{0.22\textwidth}p{0.12\textwidth}p{0.46\textwidth}p{0.14\textwidth}}
\toprule
\textbf{Name} & \textbf{Unit} & \textbf{Description} & \textbf{Default} \\
\midrule
\endhead
px & 1 & X-component of polarisation vector (can be negative) & 0 \\
py & 1 & Y-component of polarisation vector (can be negative) & 0 \\
pz & 1 & Z-component of polarisation vector (can be negative) & 0 \\
randomOn & 1 & Generate random values if randomOn!=0. & 0 \\
normalize & 1 & Normalize the polarization vector to unity length. & 0 \\
\bottomrule
\end{longtable}

\subsection*{Links}
\begin{itemize}
  \item \href{run:/home/willend/willend-McCode/mcstas-comps/optics/Set_pol.comp}{Source code} for \texttt{Set\_pol.comp}.
\end{itemize}
\IfFileExists{Set_pol_static.tex}{\input{Set_pol_static.tex}}{}